% Generated draft: requires mathematical and source review.
\documentclass[11pt]{article}
\usepackage[margin=25mm]{geometry}
\usepackage{fontspec}
\setmainfont{Latin Modern Roman}
\usepackage{amsmath,amssymb,mathtools,unicode-math}
\setmathfont{Latin Modern Math}
\usepackage{enumitem,needspace,xurl,hyperref}
\setlength{\parindent}{0pt}
\setlength{\parskip}{5pt}
\setlength{\emergencystretch}{3em}
\newcommand{\N}{\mathbb{N}}
\newcommand{\Z}{\mathbb{Z}}
\newcommand{\Q}{\mathbb{Q}}
\newcommand{\R}{\mathbb{R}}
\newcommand{\C}{\mathbb{C}}
\newcommand{\F}{\mathbb{F}}
\newcommand{\A}{\mathbb{A}}
\newcommand{\PP}{\mathbb{P}}
\newcommand{\E}{\mathbb{E}}
\newcommand{\eps}{\varepsilon}
\newcommand{\dd}{\,\mathrm{d}}
\begin{document}
\section*{UnsolvedMath ID: 30004267}
\textbf{Source problem code:} catalogue entry
\subsection{Definitions and mathematical statement}
The supplied record names a group \(G=\operatorname{Aut}_{Q}(M)\), a space \(X=\operatorname{Fl}^{mf}(V)\), and a natural action \(G\curvearrowright X\). It does not define \(Q\), \(M\), \(V\), the superscript \(mf\), or the action. These symbols are therefore retained as source notation, without assigning unsupported hypotheses.

For a specified action, the orbit of \(x\in X\) and the orbit set are
\[
G\cdot x=\{g\cdot x:g\in G\},\qquad G\backslash X=\{G\cdot x:x\in X\}.
\]
The requested problem is to determine whether
\[
\lvert G\backslash X\rvert<\infty,
\]
to give an explicit parametrization \(\lambda\mapsto\mathcal O_\lambda\) that is a bijection from a described parameter set \(\Lambda\) onto \(G\backslash X\), and to determine exactly which pairs \(\lambda,\mu\in\Lambda\) satisfy
\[
\mathcal O_\lambda\subseteq\overline{\mathcal O_\mu}^{\,X}.
\]
Here the closure must be taken in the intended topology on \(X\), which the supplied record does not specify. The underlying field, permitted objects, flag conditions, and scope of quantification over \(Q,M,V\) are also absent. Consequently, this is a precise description of the requested orbit questions conditional on the missing setup, rather than a fully specified mathematical problem.
\subsection{Short English statement}
For the action denoted by \(\operatorname{Aut}_{Q}(M)\curvearrowright\operatorname{Fl}^{mf}(V)\), determine whether there are finitely many orbits, describe parameters distinguishing them, and characterize inclusion of one orbit in the closure of another. The supplied record omits the definitions and hypotheses needed to specify the action, flags, and closure topology.
\subsection{Sources}
\begin{itemize}
\item Supplied local UnsolvedMath registry record, permanent ID \(30004267\), exact source problem code \texttt{catalogue entry}, titled ``Orbit Structure of Maximal Flat Flag Degenerations.'' No individual proposer is supplied; the proposed year is recorded as 2019.
\item Supplied report citation: \emph{Mini-Workshop: Degeneration Techniques in Representation Theory} (2020), Oberwolfach Reports, pp.~2869--2909, DOI \texttt{10.4171/owr/2019/46}. \url{https://doi.org/10.4171/owr/2019/46}.
\item Supplied literature source, labeled ``OWR 2019 46'': \url{https://ems.press/content/serial-article-files/46825}.
\item Literature assessment dated exactly \texttt{2026-08-22}: the supplied curation labels the problem open, states that the report explicitly requests the closure relation for a proposed stratification, and reports that no general closure theorem resolving it was found. This is supplied dated curation, not an independently verified literature assessment or a proof of current openness.
\item The supplied record labels its statement \texttt{corrected\_verified} and reports manual semantic, typographic, and LaTeX review asserting self-containment. That assertion is preserved as source metadata; the missing mathematical setup remains unresolved in this draft. Neither linked source has been independently consulted here.
\end{itemize}
\end{document}
